\documentclass[12pt]{article}
\usepackage[utf8]{inputenc}
\usepackage[T1]{fontenc}
\usepackage{hyperref}
\title{Pinjarup Language Resources for K-12 Curriculum Development}
\author{Compiled by Claude Sonnet 5.5 (claude-sonnet-5-5) for Andr\'as Kornai}
\date{}
\begin{document}
\maketitle

\section{Introduction}
Pinjarup (ISO 639-3 pnj) is an extinct Aboriginal language of south-western Western Australia, named after the Pinjarra area. It is usually regarded as a variety of Noongar, though that identification is not completely secure. No Pinjarup-specific resources were found; the material below concerns Noongar generally and may or may not cover Pinjarup.

\section{Primary Educational Resources}
No Pinjarup teaching materials were found.

\section{Linguistic and Reference Materials}
\begin{itemize}
\item Noongar language overview, including Pinjarup among the dialects: \url{https://en.wikipedia.org/wiki/Noongar_language}
\item A Nyoongar wordlist from the south-west of Western Australia, compiled and edited by Peter Bindon and Ross Chadwick (catalogue record; whether it contains Pinjarup entries was not confirmed): \url{https://link.stirling.wa.gov.au/resource/k8mq-i6z1xE}
\item Noongar words list, Noongar Dictionary (Maya Keniny): \url{https://www.mayakeniny.com/language-list}
\end{itemize}

\section{Religious and Community Resources}
None found specific to Pinjarup.

\section{Digital and Audio Resources}
None found specific to Pinjarup.

\section{Curriculum Development Context}
With no speakers, a heritage and revitalisation curriculum would have to draw on Noongar materials and any surviving historical word lists, and would need to involve the relevant Noongar community, whom the searches did not identify. Little exists specifically for Pinjarup.

\end{document}